\documentclass[10pt,a4paper,landscape]{article}
\usepackage[portuguese]{babel}
\usepackage[utf8]{inputenc}
\usepackage[T1]{fontenc}
\usepackage[a4paper,landscape,margin=0.34cm]{geometry}
\usepackage{amsmath,amssymb}
\pagestyle{empty}
\setlength{\parindent}{0pt}
\setlength{\parskip}{0pt}
\setlength{\abovedisplayskip}{0.45pt}
\setlength{\belowdisplayskip}{0.45pt}
\setlength{\abovedisplayshortskip}{0.2pt}
\setlength{\belowdisplayshortskip}{0.2pt}
\newcommand{\secao}[1]{\vspace{0.8mm}\textbf{\normalsize #1}\par}
\newcommand{\subsecao}[1]{\vspace{0.35mm}\textit{#1}\par}
\newcommand{\colstart}{\begin{minipage}[t]{0.244\textwidth}\vspace{0pt}\fontsize{8.1}{8.55}\selectfont}
\newcommand{\colend}{\end{minipage}}
\begin{document}
\fontsize{10.5}{11}\selectfont
\textbf{Formulário -- Segunda Parte}\par
\vspace{0.45mm}
\noindent
\colstart
\secao{Forças e Leis de Newton}
\[
\vec F_{\mathrm{res}}=\sum_{i=1}^{n}\vec F_i
=\vec F_1+\vec F_2+\cdots+\vec F_n
\]
\[
\vec F_{\mathrm{res}}=m\vec a,\qquad
F_{\mathrm{res},x}=ma_x,\quad F_{\mathrm{res},y}=ma_y
\]
\[
1\,\mathrm{N}=1\,\mathrm{kg}\cdot \mathrm{m/s^2}
\]
\subsecao{Força gravítica e peso}
\[
\vec F_g=m\vec g,\qquad P=|\vec F_g|=mg
\]
\subsecao{Atrito}
\[
|\vec f|=\mu F_N
\]

\secao{Movimento Circular}
\subsecao{Descrição angular}
\[
\theta=\frac{s}{r},\qquad
\Delta\theta=\theta_2-\theta_1=\frac{\Delta s}{r}
\]
\[
1\ \text{volta}=2\pi\ \mathrm{rad}=360^\circ
\]
\[
x(t)=r\cos\theta(t),\qquad y(t)=r\sin\theta(t)
\]
\[
\omega_{\text{méd}}=\frac{\Delta\theta}{\Delta t},
\qquad \omega=\frac{d\theta}{dt}
\]
\[
\alpha_{\text{méd}}=\frac{\Delta\omega}{\Delta t},
\qquad
\alpha=\frac{d\omega}{dt}=\frac{d^2\theta}{dt^2}
\]
\subsecao{Relações lineares e angulares}
\[
s=r\theta,\qquad v=r\omega
\]
\[
a_t=r\alpha,\qquad a_r=\frac{v^2}{r}=\omega^2r
\]
\[
\sum F_r=ma_r=m\frac{v^2}{r}=m\omega^2r
\]
\subsecao{MCU}
\[
\omega=\text{constante},\qquad \alpha=0
\]
\[
T=\frac{2\pi}{\omega}=\frac{2\pi r}{v},
\qquad f=\frac{1}{T}
\]
\[
\omega=\frac{2\pi}{T}=2\pi f,\qquad
a_c=\frac{v^2}{r}=\omega^2r
\]
\colend\hfill
\colstart
\subsecao{MCUV: $\alpha$ constante}
\[
\omega(t)=\omega(0)+\alpha t
\]
\[
\theta(t)=\theta(0)+\omega(0)t+\frac{1}{2}\alpha t^2
\]
\[
\omega^2(t)=\omega^2(0)+2\alpha[\theta(t)-\theta(0)]
\]
\[
\theta(t)=\theta(0)+\frac{1}{2}[\omega(0)+\omega(t)]t
\]
\[
\omega_{\text{méd}}=\frac{1}{2}[\omega(0)+\omega(t)]
\]

\secao{Energia Cinética e Trabalho}
\subsecao{Energia cinética}
\[
K=\frac{1}{2}mv^2
\]
\[
1\,\mathrm{J}=1\,\mathrm{kg}\,\frac{\mathrm{m^2}}{\mathrm{s^2}}
=1\,\mathrm{N}\cdot\mathrm{m}
\]
\subsecao{Trabalho de uma força constante}
\[
W=Fd\cos\phi=\vec F\cdot\vec d=F_xd_x+F_yd_y
\]
\[
F_{\parallel}=F\cos\phi,\qquad W=F_{\parallel}d
\]
\[
\vec v\cdot\vec w=vw\cos\phi
\]
\[
F_\perp=F\sin\phi\quad\Longrightarrow\quad W_\perp=0
\]
\[
0^\circ\leq\phi<90^\circ\Rightarrow W>0
\]
\[
90^\circ<\phi\leq180^\circ\Rightarrow W<0
\]
\subsecao{Teorema trabalho--energia}
\[
\Delta K=K_f-K_i=W_{\mathrm{tot}}
\]
\subsecao{Força elástica: lei de Hooke}
\[
\vec F_k=-k\vec x,\qquad F_k=k|x|
\]
\[
W_k=\frac{1}{2}kx_i^2-\frac{1}{2}kx_f^2
\]
Se $x_i=0$ e $x_f=d$:
\[
W_k=-\frac{1}{2}kd^2
\]
\colend\hfill
\colstart
\secao{Energia Potencial e Conservação}
\subsecao{Separação do trabalho}
\[
W_{\mathrm{tot}}=W_{\mathrm{cons}}+W_{\mathrm{ncons}}
\]
\[
\Delta U=U_f-U_i=-W_{\mathrm{cons}}
\]
\[
\Delta K+\Delta U=W_{\mathrm{ncons}}
\]
\subsecao{Energia potencial gravitacional}
\[
\Delta U_g=mg(y_f-y_i)=mg\Delta y
\]
Escolhendo $U=0$ em $y=0$:
\[
U_g=mgy
\]
\subsecao{Energia potencial elástica}
\[
U_k=\frac{1}{2}kx^2
\]
\[
\Delta U_k=\frac{1}{2}kx_f^2-\frac{1}{2}kx_i^2
\]
\subsecao{Energia mecânica}
\[
E_{\mathrm{mec}}=K+U
\]
Se $W_{\mathrm{ncons}}=0$:
\[
\Delta E_{\mathrm{mec}}=\Delta K+\Delta U=0
\]
\[
K_i+U_i=K_f+U_f
\]
Se $W_{\mathrm{ncons}}\neq0$:
\[
W_{\mathrm{ncons}}=\Delta E_{\mathrm{mec}}
=\Delta K+\Delta U
\]
\subsecao{Expressões úteis}
\[
K=\frac{1}{2}mv^2,\qquad U_g=mgy,\qquad U_k=\frac{1}{2}kx^2
\]

\secao{Centro de Massa}
Para $M=m_1+m_2+\cdots+m_n$:
\[
\vec r_{\mathrm{cm}}
=
\frac{m_1\vec r_1+m_2\vec r_2+\cdots+m_n\vec r_n}{M}
\]
\subsecao{Movimento do centro de massa}
\[
\vec v_{\mathrm{cm}}
=
\frac{d\vec r_{\mathrm{cm}}}{dt}
=
\frac{m_1\vec v_1+\cdots+m_n\vec v_n}{M}
\]
\[
\vec a_{\mathrm{cm}}
=
\frac{d\vec v_{\mathrm{cm}}}{dt}
=
\frac{m_1\vec a_1+\cdots+m_n\vec a_n}{M}
\]
\[
M\vec a_{\mathrm{cm}}=\vec F_{\mathrm{ext}}
\]
\colend\hfill
\colstart
\secao{Momento Linear e Impulso}
\subsecao{Partícula}
\[
\vec p=m\vec v,\qquad [\vec p]=\mathrm{kg}\,\mathrm{m/s}
\]
\subsecao{Sistema de partículas}
\[
\vec P=\vec p_1+\vec p_2+\cdots+\vec p_n
\]
\[
\vec P=m_1\vec v_1+m_2\vec v_2+\cdots+m_n\vec v_n
\]
\[
\vec P=M\vec v_{\mathrm{cm}}
\]
\subsecao{Segunda lei para um sistema}
\[
\frac{d\vec P}{dt}=\vec F_{\mathrm{ext}}
\]
\subsecao{Conservação do momento linear}
Se $\vec F_{\mathrm{ext}}=\vec 0$:
\[
\vec P=\text{constante}\quad\Longrightarrow\quad \vec P_i=\vec P_f
\]
\subsecao{Impulso}
\[
\Delta\vec p=\vec p_f-\vec p_i,\qquad \vec J=\Delta\vec p
\]
\[
F=\frac{\Delta p}{\Delta t},\qquad \Delta p=|\Delta\vec p|
\]

\secao{Colisões}
\subsecao{Sistema fechado e isolado}
\[
\vec P_i=\vec P_f
\]

\subsecao{Colisão elástica}
\[
K_i=K_f
\]
\subsecao{Colisão inelástica}
\[
\vec P_i=\vec P_f,\qquad K_i\neq K_f
\]
Normalmente:
\[
K_f<K_i
\]
\subsecao{Perfeitamente inelástica}
\[
\vec P_i=\vec P_f,\qquad K_f<K_i
\]
Os corpos ficam juntos depois da colisão.
\colend
\end{document}
